\documentclass[10pt,a4paper]{amsart}
\usepackage{amsmath,amssymb,amsthm,graphicx}
\usepackage{verbatim}
\newenvironment{reference}{\comment}{\endcomment}
%\newenvironment{reference}{}{}
\newenvironment{slogan}{\comment}{\endcomment}
\newenvironment{history}{\comment}{\endcomment}

% For commutative diagrams we use Xy-pic
\usepackage[all]{xy}

% We use 2cell for 2-commutative diagrams.
\xyoption{2cell}
% Theorem environments.
%
\theoremstyle{plain}
\newtheorem{theorem}[subsection]{Theorem}
\newtheorem{proposition}[subsection]{Proposition}
\newtheorem{lemma}[subsection]{Lemma}

\theoremstyle{definition}
\newtheorem{definition}[subsection]{Definition}
\newtheorem{example}[subsection]{Example}
\newtheorem{exercise}[subsection]{Exercise}
\newtheorem{situation}[subsection]{Situation}

\theoremstyle{remark}
\newtheorem{remark}[subsection]{Remark}
\newtheorem{remarks}[subsection]{Remarks}

\numberwithin{equation}{subsection}

% Macros
%
\def\lim{\mathop{\mathrm{lim}}\nolimits}
\def\colim{\mathop{\mathrm{colim}}\nolimits}
\def\Spec{\mathop{\mathrm{Spec}}}
\def\Hom{\mathop{\mathrm{Hom}}\nolimits}
\def\Ext{\mathop{\mathrm{Ext}}\nolimits}
\def\SheafHom{\mathop{\mathcal{H}\!\mathit{om}}\nolimits}
\def\SheafExt{\mathop{\mathcal{E}\!\mathit{xt}}\nolimits}
\def\Sch{\mathit{Sch}}
\def\Mor{\mathop{\mathrm{Mor}}\nolimits}
\def\Ob{\mathop{\mathrm{Ob}}\nolimits}
\def\Sh{\mathop{\mathit{Sh}}\nolimits}
\def\NL{\mathop{N\!L}\nolimits}
\def\CH{\mathop{\mathrm{CH}}\nolimits}
\def\proetale{{pro\text{-}\acute{e}tale}}
\def\etale{{\acute{e}tale}}
\def\QCoh{\mathit{QCoh}}
\def\Ker{\mathop{\mathrm{Ker}}}
\def\Im{\mathop{\mathrm{Im}}}
\def\Coker{\mathop{\mathrm{Coker}}}
\def\Coim{\mathop{\mathrm{Coim}}}

% Boxtimes
%
\DeclareMathSymbol{\boxtimes}{\mathbin}{AMSa}{"02}

%
% Macros for moduli stacks/spaces
%
\def\QCohstack{\mathcal{QC}\!\mathit{oh}}
\def\Cohstack{\mathcal{C}\!\mathit{oh}}
\def\Spacesstack{\mathcal{S}\!\mathit{paces}}
\def\Quotfunctor{\mathrm{Quot}}
\def\Hilbfunctor{\mathrm{Hilb}}
\def\Curvesstack{\mathcal{C}\!\mathit{urves}}
\def\Polarizedstack{\mathcal{P}\!\mathit{olarized}}
\def\Complexesstack{\mathcal{C}\!\mathit{omplexes}}
% \Pic is the operator that assigns to X its picard group, usage \Pic(X)
% \Picardstack_{X/B} denotes the Picard stack of X over B
% \Picardfunctor_{X/B} denotes the Picard functor of X over B
\def\Pic{\mathop{\mathrm{Pic}}\nolimits}
\def\Picardstack{\mathcal{P}\!\mathit{ic}}
\def\Picardfunctor{\mathrm{Pic}}
\def\Deformationcategory{\mathcal{D}\!\mathit{ef}}
\usepackage{cleveref}
\crefname{originalequation}{original equation}{original equations}
\Crefname{originalequation}{Original equation}{Original equations}
\makeatletter
\newlabel{equation-completion-functor}{{0880}{source}{}{sourcepointer-1}{}}
\newlabel{equation-completion-functor@cref}{{[originalequation][0][]0880}{[1][0][]source}{}{}{}}
\makeatother
\crefname{originallemma}{original lemma}{original lemmas}
\Crefname{originallemma}{Original lemma}{Original lemmas}
\makeatletter
\newlabel{formal-spaces-lemma-reduction-completion}{{0GB9}{source}{}{sourcepointer-2}{}}
\newlabel{formal-spaces-lemma-reduction-completion@cref}{{[originallemma][0][]0GB9}{[1][0][]source}{}{}{}}
\makeatother
\crefname{originallemma}{original lemma}{original lemmas}
\Crefname{originallemma}{Original lemma}{Original lemmas}
\makeatletter
\newlabel{formal-spaces-lemma-reduction-surjective}{{0GB8}{source}{}{sourcepointer-3}{}}
\newlabel{formal-spaces-lemma-reduction-surjective@cref}{{[originallemma][0][]0GB8}{[1][0][]source}{}{}{}}
\makeatother
\crefname{originalreference}{original reference}{original references}
\Crefname{originalreference}{Original reference}{Original references}
\makeatletter
\newlabel{lemma-algebraize-morphism-rig-etale}{{0AR3}{source}{}{sourcepointer-4}{}}
\newlabel{lemma-algebraize-morphism-rig-etale@cref}{{[originalreference][0][]0AR3}{[1][0][]source}{}{}{}}
\makeatother
\crefname{originalreference}{original reference}{original references}
\Crefname{originalreference}{Original reference}{Original references}
\makeatletter
\newlabel{lemma-algebraize-rig-etale-affine}{{0AR4}{source}{}{sourcepointer-5}{}}
\newlabel{lemma-algebraize-rig-etale-affine@cref}{{[originalreference][0][]0AR4}{[1][0][]source}{}{}{}}
\makeatother
\crefname{originalreference}{original reference}{original references}
\Crefname{originalreference}{Original reference}{Original references}
\makeatletter
\newlabel{lemma-faithful}{{0240}{source}{}{sourcepointer-6}{}}
\newlabel{lemma-faithful@cref}{{[originalreference][0][]0240}{[1][0][]source}{}{}{}}
\makeatother
\title{Completion along a closed subspace induces a bijection on morphisms into a modification which is an isomorphism over the open}
\author{The Stacks Project Authors}
\date{stacks-project:restricted.tex:local-snapshot-20261012; GNU FDL 1.2 or later}
\begin{document}
\maketitle
\noindent\textit{Editorial note.} All delivered mathematical text is the author's own, copied from the pinned Stacks Project chapter file restricted.tex; no statement or proof step was rewritten, reordered or removed.\par\medskip
\noindent\textit{Editorial note.} The counted proof is the author's complete argument: an injectivity step, an etale covering gluing step, and a two-stage reduction to the affine case closed by the rig-etale algebraization lemma, all carried out inside the delivered span\par\medskip
\noindent\textit{Editorial note.} The label lemma-completion-functor-fully-faithful names the counted statement, and the proof environment immediately following it in the source is delivered verbatim as its proof.\par\medskip
\noindent\textit{Editorial note.} The counted statement is carried by the author's own lemma environment in the chapter source.\par\medskip
\noindent\textit{Editorial note.} Revision r002: regenerated from the same spec after the r001 assembly crashed during packaging (stale empty job directory retained as evidence).\par\medskip
\noindent\textit{Allowed standard prerequisites (editorial).}\begin{itemize}
\item The completion construction Y maps to Y\_\{/T\} for algebraic spaces over the pair (X, T), as set up in the surrounding section of the source
\item The fppf sheaf property of the algebraic space X' and descent of morphisms along etale coverings
\item Formal algebraic spaces: the reduction of a completion is the reduced induced scheme on the closed subspace
\end{itemize}\medskip
\ifdefined\hypertarget\hypertarget{sourcepointer-1}{}\fi
\noindent\textit{Source reference.} Original equation 0880: sources/stacks-restricted/pinned-original-source.tex, source lines 6766--6766. Prerequisite pointer invoked at proof line(s) 6840; it names a standard fact of the same chapter that the counted argument uses, and every step of the counted proof built on it is retained verbatim in the delivered body.\par\medskip
\ifdefined\hypertarget\hypertarget{sourcepointer-2}{}\fi
\noindent\textit{Source reference.} Original lemma 0GB9: sources/stacks-chow/tags.txt, source lines 20364--20364. Prerequisite pointer invoked at proof line(s) 6896; it names a standard fact of a sibling chapter that the counted argument uses, and every step of the counted proof built on it is retained verbatim in the delivered body.\par\medskip
\ifdefined\hypertarget\hypertarget{sourcepointer-3}{}\fi
\noindent\textit{Source reference.} Original lemma 0GB8: sources/stacks-chow/tags.txt, source lines 20363--20363. Prerequisite pointer invoked at proof line(s) 6914; it names a standard fact of a sibling chapter that the counted argument uses, and every step of the counted proof built on it is retained verbatim in the delivered body.\par\medskip
\ifdefined\hypertarget\hypertarget{sourcepointer-4}{}\fi
\noindent\textit{Source reference.} Original reference 0AR3: sources/stacks-restricted/pinned-original-source.tex, source lines 6425--6425. Prerequisite pointer invoked at proof line(s) 6856; it names a standard fact of the same chapter that the counted argument uses, and every step of the counted proof built on it is retained verbatim in the delivered body.\par\medskip
\ifdefined\hypertarget\hypertarget{sourcepointer-5}{}\fi
\noindent\textit{Source reference.} Original reference 0AR4: sources/stacks-restricted/pinned-original-source.tex, source lines 6406--6406. Prerequisite pointer invoked at proof line(s) 6915; it names a standard fact of the same chapter that the counted argument uses, and every step of the counted proof built on it is retained verbatim in the delivered body.\par\medskip
\ifdefined\hypertarget\hypertarget{sourcepointer-6}{}\fi
\noindent\textit{Source reference.} Original reference 0240: sources/stacks-restricted/pinned-original-source.tex, source lines 6502--6502. Prerequisite pointer invoked at proof line(s) 6848; it names a standard fact of the same chapter that the counted argument uses, and every step of the counted proof built on it is retained verbatim in the delivered body.\par\medskip
\begin{lemma}
\label{lemma-completion-functor-fully-faithful}
In the situation above. Let $f : X' \to X$ be a morphism of algebraic spaces
which is locally of finite type and an isomorphism over $U$. Let
$g : Y \to X$ be a morphism with $Y$ locally Noetherian. Then completion
defines a bijection
$$
\Mor_X(Y, X') \longrightarrow \Mor_{X_{/T}}(Y_{/T}, X'_{/T})
$$
In particular, the functor (\ref{equation-completion-functor}) is
fully faithful.
\end{lemma}

\begin{proof}
Let $a, b : Y \to X'$ be morphisms over $X$ such that
$a_{/T} = b_{/T}$. Then we see that $a$ and $b$ agree over the open
subspace $g^{-1}U$ and after completion along $g^{-1}T$.
Hence $a = b$ by Lemma \ref{lemma-faithful}.
In other words, the completion map is always injective.

\medskip\noindent
Let $\alpha : Y_{/T} \to X'_{/T}$ be a morphism of formal algebraic spaces
over $X_{/T}$. We have to prove there exists a morphism $a : Y \to X'$
over $X$ such that $\alpha = a_{/T}$. The proof proceeds by a standard
but cumbersome reduction to the affine case and then applying
Lemma \ref{lemma-algebraize-morphism-rig-etale}.

\medskip\noindent
Let $\{h_i : Y_i \to Y\}$ be an \'etale covering of algebraic spaces.
If we can find for each $i$ a morphism $a_i : Y_i \to X'$ over $X$
whose completion $(a_i)_{/T} : (Y_i)_{/T} \to X'_{/T}$ is equal to
$\alpha \circ (h_i)_{/T}$, then we get a morphism $a : Y \to X'$
with $\alpha = a_{/T}$. Namely, we first observe that
$(a_i)_{/T} \circ \text{pr}_1 = (a_j)_{/T} \circ \text{pr}_2$
as morphisms $(Y_i \times_Y Y_j)_{/T} \to X'_{/T}$ by the
agreement with $\alpha$ (this uses that completion ${}_{/T}$ commutes
with fibre products). By the injectivity already proven this shows that
$a_i \circ \text{pr}_1 = a_j \circ \text{pr}_2$ as morphisms
$Y_i \times_Y Y_j \to X'$. Since $X'$ is an fppf sheaf this means that
the collection of morphisms $a_i$ descends to a morphism $a : Y \to X'$.
We have $\alpha = a_{/T}$ because $\{(a_i)_{/T} : (Y_i)_{/T} \to X'_{/T}\}$
is an \'etale covering.

\medskip\noindent
By the result of the previous paragraph, to prove existence,
we may assume that $Y$ is affine and that $g : Y \to X$ factors
as $g_1 : Y \to X_1$ and an \'etale morphism $X_1 \to X$ with $X_1$ affine.
Then we can consider $T_1 \subset |X_1|$ the inverse image of $T$
and we can set $X'_1 = X' \times_X X_1$ with projection $f_1 : X'_1 \to X_1$
and
$$
\alpha_1 = (\alpha, (g_1)_{/T_1}) :
Y_{/T_1} = Y_{/T}
\longrightarrow
X'_{/T} \times_{X_{/T}} (X_1)_{/T_1} = (X'_1)_{/T_1}
$$
We conclude that it suffices to prove the existence for $\alpha_1$
over $X_1$, in other words, we may replace $X, T, X', Y, f, g, \alpha$
by $X_1, T_1, X'_1, Y, g_1, \alpha_1$.
This reduces us to the case described in the next paragraph.

\medskip\noindent
Assume $Y$ and $X$ are affine. Recall that $(Y_{/T})_{red}$
is an affine scheme (isomorphic to the reduced induced scheme structure
on $g^{-1}T \subset Y$, see Formal Spaces, Lemma
\ref{formal-spaces-lemma-reduction-completion}).
Hence $\alpha_{red} : (Y_{/T})_{red} \to (X'_{/T})_{red}$
has quasi-compact image $E$ in $f^{-1}T$ (this is the underlying
topological space of $(X'_{/T})_{red}$ by the same lemma as above).
Thus we can find an affine scheme $V$ and an \'etale morpism
$h : V \to X'$ such that the image of $h$ contains $E$.
Choose a solid cartesian diagram
$$
\xymatrix{
Y'_{/T} \ar@{..>}[rd] \ar@{..>}[r] &
W \ar[d] \ar[r] & V_{/T} \ar[d]^{h_{/T}} \\
& Y_{/T} \ar[r]^\alpha & X'_{/T}
}
$$
By construction, the morphism $W \to Y_{/T}$ is representable by
algebraic spaces, \'etale, and surjective
(surjectivity can be seen by looking at the reductions, see
Formal Spaces, Lemma
\ref{formal-spaces-lemma-reduction-surjective}).
By Lemma \ref{lemma-algebraize-rig-etale-affine}
we can write $W = Y'_{/T}$ for $Y' \to Y$ \'etale and $Y'$ affine.
This gives the dotted arrows in the diagram.
Since $W \to Y_{/T}$ is surjective, we see that
the image of $Y' \to Y$ contains $g^{-1}T$.
Hence $\{Y' \to Y, Y \setminus g^{-1}T \to Y\}$
is an \'etale covering. As $f$ is an isomorphism over $U$ we have a
(unique) morphism $Y \setminus g^{-1}T \to X'$ over $X$
agreeing with $\alpha$ on completions
(as the completion of $Y \setminus g^{-1}T$ is empty).
Thus it suffices to prove the existence for $Y'$
which reduces us to the case studied in the next paragraph.

\medskip\noindent
By the result of the previous paragraph, we may assume that $Y$ is affine
and that $\alpha$ factors as $Y_{/T} \to V_{/T} \to X'_{/T}$
where $V$ is an affine scheme \'etale over $X'$. We may still replace $Y$ by
the members of an affine \'etale covering.
By Lemma \ref{lemma-algebraize-morphism-rig-etale}
we may find an \'etale morphism $b : Y' \to Y$ of affine schemes
which induces an isomorphism $b_{/T} : Y'_{/T} \to Y_{/T}$
and a morphism $c : Y' \to V$ such that $c_{/T} \circ b_{/T}^{-1}$
is the given morphism $Y_{/T} \to V_{/T}$. Setting
$a' : Y' \to X'$ equal to the composition of $c$ and $V \to X'$
we find that $a'_{/T} = \alpha \circ b_{/T}$, in other words, we have
existence for $Y'$ and $\alpha \circ b_{/T}$.
Then we are done by replacing considering once more
the \'etale covering $\{Y' \to Y, Y \setminus g^{-1}T \to Y\}$.
\end{proof}

\end{document}
